" 
 ملحوظة الي هيكون بين علامتين تنصيص هو الي اتعمل اما من غير تنصيص لسة
 "
"عمل اتصال بين السيرفر ""والمتصفح ""
 عملت ربط بين الباك اند وقاعدة البيانات
 ""عملت استعادة المستخدم "
 {
    انشاء حساب ,
    تسجيل دخول ,
    حذف حساب,
    تحديث حساب 
 }
 البوست{
    انشاء بوست,
    تجديث بوست ,
    حذف بوست
 }
 رفع الصور من multer
  الاخطاء المهمة
  200-299: معناها "الدنيا تمام" (Success).

400-499: معناها "فيه غلطة من عند المستخدم" (Client Error) زي باسوورد غلط.

500-599: معناها "السيرفر فيه مشكلة" (Server Error) زي الداتابيز واقعة.






لغة .md
قال Gemini
امتداد .md هو اختصار لـ Markdown، وهو ببساطة لغة "تنسيق نصي" بسيطة جداً صُممت عشان المبرمجين والتقنيين يكتبوا نصوص منسقة (فيها عناوين، صور، لينكات، أكواد) باستخدام نص عادي (Plain Text) من غير ما يحتاجوا لبرامج تقيلة زي Word أو كتابة كود HTML معقد.

بما إنك شغال على GitHub وبتبني سيستم بـ Fastify، فالامتداد ده هو "اللغة الرسمية" لوصف مشاريعك.

ليه الامتداد ده مهم جداً ليك كمبرمج؟
ملف الـ README.md: هو الواجهة الحقيقية لأي مشروع ليك على GitHub. من غيره، محدش هيفهم السيستم بتاعك بيعمل إيه.

السهولة: بدل ما تكتب كود HTML عشان تعمل عنوان كبير <h1>Title</h1> بمجرد ما تحط # Title في ملف الـ .md هيتحول لعنوان.

دعم الأكواد: بيسمحلك تعرض الكود بتاعك (زي كود الـ Fastify أو Python) بشكل ملون ومنظم (Syntax Highlighting) جوه ملف النص.

أشهر العلامات اللي هتحتاجها في الـ Markdown:
النتيجة المتوقعة	الطريقة في ملف الـ .md
عنوان رئيسي	# اسم العنوان
نص عريض (Bold)	**كلمة مهمة**
رابط (Link)	[اسم الرابط](URL)
قائمة نقاط	* عنصر 1
بلوك كود	استخدام الثلاث علامات ( ``` ) قبل وبعد الكو